\section{Perfect-square conformal sector}
\label{sec:model}

We study the four-derivative scalar theory with Euclidean action
\begin{equation}
  S_{\mathrm{PS}}[\sigma]
  = \frac{1}{2}\int \dd^4x\,
  \left[\Box\sigma + \lambda\,(\partial_\mu\sigma)(\partial^\mu\sigma)\right]^2.
  \label{eq:ps-action}
\end{equation}
The overall Euclidean sign is chosen so that the action is nonnegative.  Expanding the square fixes the relative coefficients of the kinetic, cubic, and quartic derivative interactions.  Thus the model occupies a distinguished trajectory in the larger space of shift-symmetric four-derivative scalar theories rather than representing a generic derivative expansion.

The field $\sigma$ may be interpreted as the local scale factor in a conformally flat metric,
\begin{equation}
  g_{\mu\nu}=e^{2\lambda\sigma}\delta_{\mu\nu}.
  \label{eq:conformal-metric}
\end{equation}
On this restricted configuration space, the curvature-squared action reduces to the perfect-square structure in eq.~\eqref{eq:ps-action}.  This relation supplies a gravitational Ward identity that protects the perfect-square coupling relation under renormalization~\cite{Anderson:2026}.

The model does not contain propagating tensor gravitons.  It consequently cannot decide whether a proposed resolution of indefinite-norm states extends to the Weyl or spin-two sector of full quadratic gravity.  It can nevertheless test a necessary part of that program: whether an interacting, asymptotically free, gravity-derived four-derivative sector is nonperturbatively well defined.
